%!TEX root = ./main.tex

\section[Background]{Background: The Disaggregated RAN}

% TALK MAPPING: this section is slides 2--3 of the authors' SIGCOMM'25 deck,
% transcript 00:48--04:14. Two ideas only:
%   (1) the 5G RAN is now three separable, softwarized boxes with open interfaces;
%   (2) that openness has NOT produced the features it promised.
% COLOR CONVENTION for the whole deck -- do not break it:
%   ranblue    = RAN network functions (CU/DU/RU)
%   softgreen  = anything spliced into the path (the RIC, a RANBooster middlebox)
%   talkred    = the one claim the slide is really making
%   coreorange = the cellular core, and the cost side of a tradeoff
% LAYOUT RULES: a frame title must fit one line in the title bar (~44 characters).

% Teaching note (90 s): build the chain left to right and say what the chain is FOR --
% packets in from the core, radio signals out of the antenna. Do not name E2/O1 yet.
\begin{frame}{Recent trends in cellular RAN technologies}
  \vspace{0.2em}
  {\small\textbf{5G RAN characteristics:}\par}
  \vspace{0.15em}
  {\small
   \begin{itemize}\itemsep1pt
     \item Standardized and open interfaces
     \item Softwarized on commodity hardware \expand{x86, ARM, GPUs}
   \end{itemize}\par}

  \vspace{0.4em}
  \centering
  \begin{tikzpicture}[scale=0.9,transform shape]
    \node[cloud,draw=coreorange,fill=orange!8,cloud puffs=11,cloud ignores aspect,
          minimum width=1.9cm,minimum height=1.0cm,align=center,
          font=\scriptsize\bfseries] (core) at (0,0) {Cellular\\core};
    \node[netnode,minimum width=1.5cm] (cu) at (3.0,0) {CU};
    \node[netnode,minimum width=1.5cm] (du) at (6.0,0) {DU};
    \runode{ru}{(8.7,0)}{RU}

    \draw[wire] (core.east) -- (cu.west);
    \node[font=\tiny,text=softgray] at (1.8,0.32) {NG-C};
    \node[font=\tiny,text=softgray] at (1.8,-0.32) {NG-U};
    \draw[wire] (cu) -- node[font=\tiny,text=softgray,above]{F1} (du);
    \draw[wire] (du) -- node[font=\tiny,text=softgray,above]{Fronthaul} (ru);

    \begin{scope}[on background layer]
      \node[draw=softgray,dashed,inner sep=8pt,fit=(cu)(du)(ru)] (ranbox) {};
    \end{scope}
    \node[font=\scriptsize\bfseries,text=softgray,below=1pt of ranbox]
      {Radio Access Network};

    % Overlay 2: the RIC, and the applications it was supposed to host.
    \node[ricnode,visible on=<2->] (ric) at (4.5,1.9) {RIC};
    \draw[wire,softgreen,visible on=<2->] (ric.south) -- (4.5,1.15);
    \draw[wire,softgreen,visible on=<2->] (3.0,1.15) -- (6.0,1.15);
    \draw[wire,softgreen,visible on=<2->] (3.0,1.15) -- (cu.north);
    \draw[wire,softgreen,visible on=<2->] (6.0,1.15) -- (du.north);
    \node[font=\tiny,text=softgray,visible on=<2->] at (6.3,1.42)
      {E2, O1 (KPIs, control)};
    \node[draw=ranblue,fill=blue!14,rounded corners=2pt,minimum width=1.0cm,
          minimum height=0.5cm,font=\scriptsize,visible on=<2->]
          (apps) at (3.45,3.05) {Apps};
    \node[font=\scriptsize,align=left,visible on=<2->] (applist) at (5.6,3.05)
      {$\cdot$~Interference\\[-1pt]$\cdot$~Mobility\\[-1pt]$\cdot$~Security};
    \begin{scope}[on background layer]
      \node[draw=softgray,dashed,inner sep=4pt,fit=(apps)(applist),
            visible on=<2->] {};
    \end{scope}
  \end{tikzpicture}

  \glossline{CU / DU / RU = Centralized / Distributed / Radio Unit \quad
    RIC = RAN Intelligent Controller}
\end{frame}

% Teaching note (120 s): this is the slide the whole paper hangs on. Read the plus
% line, then the minus line, then ask the room the question the authors ask -- "so
% build a RIC app?" -- and only then show the trade-press headline. The joke lands
% because the RIC was the official answer.
\begin{frame}[t]{Opportunities and challenges for new players}
  \vspace{0.2em}
  {\small Disaggregation, virtualization and open interfaces:\par}
  \vspace{0.3em}
  {\small\color{softgreen}\textbf{+}~~Opportunity to open the closed cellular
    ecosystem\par}
  \vspace{0.2em}
  {\small\color{talkred}\textbf{--}~~Cannot match the performance of integrated
    end-to-end solutions\par}

  \vspace{0.5em}
  \centering
  \begin{tikzpicture}[scale=0.9,transform shape]
    \foreach \i/\t in {0/{Interference\\mitigation?},
                       1/{Mobility\\management?},
                       2/{Radio and spectrum\\sharing?}}{
      \node[draw=ranblue!35,fill=blue!5,rounded corners=2pt,minimum width=2.5cm,
            minimum height=1.05cm,align=center,font=\scriptsize,text=talkred]
            at (\i*2.9,0) {\t};
    }
    \node[font=\Large,text=talkred] at (8.5,0) {$\cdots$};
  \end{tikzpicture}

  \vspace{0.5em}
  \only<2->{%
    \begin{tcolorbox}[colback=white,colframe=black,boxrule=0.7pt,arc=1pt,
      left=5pt,right=5pt,top=3pt,bottom=3pt,width=0.72\textwidth]
      \centering\normalsize\textbf{Build RIC apps to enable advanced features?}
    \end{tcolorbox}}

  \vspace{0.35em}
  \only<3->{%
    {\small\textbf{``Open RAN xApps look stillborn''}\par}
    \vspace{0.1em}
    {\scriptsize\color{softgray}The lack of support from Ericsson and Samsung for
      the near-real-time RIC has left xApps with nowhere to go.\par}
    \source{https://www.lightreading.com/open-ran}{Iain Morris, Light Reading, December 5, 2024 --- the headline the authors put on this slide}}
\end{frame}
